\frfilename{mt222.tex}
\versiondate{20.11.03/18.10.04}
\copyrightdate{2004}


%$(\DiniD^+f)(x)=(\DiniD^-f)(x)
%=(\Dinid^+f)(x)=(\Dinid^-f)(x)$

\def\chaptername{Teorema Dasar Kalkulus}
\def\sectionname{Mendiferensialkan integral tak tentu}
     
\newsection{222}
     
Sekarang saya sampai pada pertanyaan pertama dari dua pertanyaan yang
disebutkan dalam pengantar bab ini:  jika $f$ merupakan fungsi yang
terintegralkan pada $[a,b]$, berapakah $\bover{d}{dx}\int_a^xf$? Ternyata
turunan ini ada dan sama dengan $f$ hampir di mana-mana (222E). Argumennya
didasarkan pada suatu sifat mencolok dari fungsi monoton: fungsi-fungsi
tersebut terdiferensialkan hampir di mana-mana (222A), dan kita dapat
membatasi integral turunannya (222C).
     
\leader{222A}{Teorema} Misalkan $I\subseteq\Bbb R$ suatu interval dan
$f:I\to\Bbb R$ suatu fungsi monoton. Maka $f$ terdiferensialkan hampir
di mana-mana dalam $I$.
     
\cmmnt{\medskip
     
\noindent{\bf Catatan} Jika saya tampak membicarakan suatu ukuran pada
$\Bbb R$ tanpa menyebut namanya, seperti di sini, yang saya maksud ialah
ukuran Lebesgue.
}%end of comment
     
\proof{ Seperti biasa, tuliskan $\mu^*$ untuk ukuran luar Lebesgue
pada $\Bbb R$, dan $\mu$ untuk ukuran Lebesgue.
     
\medskip
     
{\bf (a)} Sebagai permulaan (sampai akhir (c) di bawah), andaikan bahwa
$f$ tak-menurun dan $I$ merupakan interval terbuka terbatas tempat $f$
terbatas; katakanlah $|f(x)|\le M$ untuk $x\in I$. Untuk setiap
subinterval tertutup $J=[a,b]$ dari $I$, tuliskan $f^*(J)$ untuk interval
terbuka $\ooint{f(a),f(b)}$. Untuk $x\in I$, tuliskan
     
\Centerline{$\DiniD f(x)=\limsup_{h\to 0}{1\over h}(f(x+h)-f(x))$,
\quad$\Dinid f(x)=\liminf_{h\to 0}{1\over h}(f(x+h)-f(x))$,}
     
\noindent dengan membolehkan nilai $\infty$ dalam kedua kasus. Maka $f$
terdiferensialkan di $x$ jika dan hanya jika
$\DiniD f(x)=\Dinid f(x)\in\Bbb R$. Karena tentu saja
$\DiniD f(x)\ge \Dinid f(x)\ge 0$, $f$ akan terdiferensialkan di $x$
jika dan hanya jika $\DiniD f(x)$ berhingga dan
$\DiniD f(x)\le \Dinid f(x)$.
     
Oleh karena itu saya harus menunjukkan bahwa himpunan-himpunan
     
\Centerline{$\{x:x\in I,\,\DiniD f(x)=\infty\}$,
\quad$\{x:x\in I,\,\DiniD f(x)>\Dinid f(x)\}$}
     
\noindent terabaikan.
     
\medskip
     
{\bf (b)} Mula-mula ambil
$A=\{x:x\in I,\,\DiniD f(x)=\infty\}$. Untuk sementara, tetapkan suatu
bilangan bulat $m\ge 1$, dan tetapkan
     
\Centerline{$A_m=\{x:x\in I,\,\DiniD f(x)>m\}\supseteq A$.}
     
\noindent Misalkan $\Cal I$ keluarga interval tertutup tak-sepele
$[a,b]\subseteq I$ sedemikian sehingga $f(b)-f(a)\ge m(b-a)$; maka
$\mu f^*(J)\ge m\mu J$ untuk setiap $J\in\Cal I$. Jika $x\in A_m$,
maka untuk setiap $\delta>0$ terdapat $h$ dengan $0<|h|\le\delta$ dan
${1\over h}(f(x+h)-f(x))>m$, sehingga
     
\Centerline{$[x,x+h]\in\Cal I$ jika $h>0$,
\quad$[x+h,x]\in\Cal I$ jika $h<0$;}
     
\noindent
dengan demikian setiap anggota $A_m$ termasuk dalam interval-interval
yang sekecil apa pun dalam $\Cal I$. Menurut teorema Vitali (221A),
terdapat suatu subkeluarga terhitung dan saling lepas
$\Cal I_0\subseteq\Cal I$
sedemikian sehingga $\mu(A_m\setminus\bigcup\Cal I_0)=0$.
Sekarang, karena $f$ tak-menurun,
$\langle f^*(J)\rangle_{J\in\Cal I_0}$ saling lepas, dan semua $f^*(J)$
termasuk dalam $[-M,M]$, sehingga
$\sum_{J\in\Cal I_0}\mu f^*(J)\le 2M$ dan
$\sum_{J\in\Cal I_0}\mu J\le 2M/m$. Karena
$A_m\setminus\bigcup\Cal I_0$ terabaikan,
     
     
\Centerline{$\mu^*A\le\mu^*A_m\le\Bover{2M}{m}$.}
     
\noindent Karena $m$ sembarang, $\mu^*A=0$ dan $A$ terabaikan.
     
\medskip
     
{\bf (c)} Sekarang tinjau
$B=\{x:x\in I,\,\DiniD f(x)>\Dinid f(x)\}$. Untuk $q$, $q'\in\Bbb Q$
dengan $0\le q<q'$, tetapkan
     
     
\Centerline{$B_{qq'}=\{x:x\in I,\,\Dinid f(x)< q,\,\DiniD f(x)>q'\}$.}
     
\noindent Untuk sementara tetapkan $q$, $q'$ seperti itu, dan tuliskan
$\gamma=\mu^*B_{qq'}$. Ambil sembarang $\epsilon>0$, dan misalkan $G$
suatu himpunan terbuka yang memuat $B_{qq'}$ sedemikian sehingga
$\mu G\le\gamma+\epsilon$ (134Fa). Misalkan $\Cal J$ keluarga interval
tertutup tak-sepele $[a,b]\subseteq I\cap G$ sedemikian sehingga
$f(b)-f(a)\le q(b-a)$; kali ini $\mu f^*(J)\le q\mu J$ untuk
$J\in\Cal J$. Maka setiap anggota $B_{qq'}$ termasuk dalam anggota-anggota
$\Cal J$ yang sekecil apa pun, sehingga terdapat suatu subkeluarga
terhitung dan saling lepas $\Cal J_0\subseteq\Cal J$ sedemikian sehingga
$B_{qq'}\setminus\bigcup\Cal J_0$ terabaikan. Misalkan $L$ himpunan
titik-titik ujung anggota $\Cal J_0$; maka $L$ merupakan gabungan terhitung
dari himpunan-himpunan beranggota dua, jadi terhitung, dan karenanya
terabaikan. Tetapkan
     
\Centerline{$C=B_{qq'}\cap\bigcup\Cal J_0\setminus L$;}
     
\noindent maka $\mu^*C=\gamma$. Misalkan $\Cal I$ keluarga interval
tertutup tak-sepele $J=[a,b]$ sedemikian sehingga (i) $J$ termasuk
dalam salah satu anggota $\Cal J_0$ (ii) $f(b)-f(a)\ge q'(b-a)$;
sekarang $\mu f^*(J)\ge q'\mu J$ untuk setiap $J\in\Cal I$. Sekali lagi,
karena setiap anggota $C$ merupakan titik interior dari suatu anggota
$\Cal J_0$, setiap titik $C$ termasuk dalam anggota-anggota $\Cal I$
yang sekecil apa pun; jadi terdapat suatu subkeluarga terhitung dan
saling lepas $\Cal I_0\subseteq\Cal I$ sedemikian sehingga
$\mu(C\setminus\bigcup\Cal I_0)=0$.
     
Seperti dalam (b) di atas,
     
\Centerline{$\gamma q'\le q'\mu(\bigcup\Cal I_0)
=\sum_{I\in\Cal I_0}q'\mu I
\le\sum_{I\in\Cal I_0}\mu f^*(I)
=\mu(\bigcup_{I\in\Cal I_0}f^*(I))$.}
     
\noindent Di sisi lain,
     
$$\eqalign{\mu(\bigcup_{J\in\Cal J_0}f^*(J))
&=\sum_{J\in\Cal J_0}\mu f^*(J)
\le q\sum_{J\in\Cal J_0}\mu J
=q\mu(\bigcup\Cal J_0)\cr
&\le q\mu(\bigcup\Cal J)
\le q\mu G
\le q(\gamma+\epsilon).\cr}$$
     
     
\noindent Tetapi
$\bigcup_{I\in\Cal I_0}f^*(I)\subseteq\bigcup_{J\in\Cal J_0}f^*(J)$,
karena setiap anggota $\Cal I_0$ termasuk dalam suatu anggota
$\Cal J_0$, sehingga $\gamma q'\le q(\gamma+\epsilon)$ dan
$\gamma\le \epsilon q/(q'-q)$. Karena $\epsilon$ sembarang, $\gamma=0$.
     
Jadi setiap $B_{qq'}$ terabaikan. Akibatnya
$B=\bigcup_{q,q'\in\Bbb Q,0\le q<q'}B_{qq'}$ terabaikan.
     
\medskip
     
{\bf (d)} Ini menangani kasus interval terbuka terbatas tempat $f$
terbatas dan tak-menurun. Masih untuk $f$ yang tak-menurun, tetapi untuk
sembarang interval $I$, perhatikan bahwa
$K=\{(q,q'):q,\,q'\in I\cap\Bbb Q,\,q<q'\}$ terhitung dan bahwa
$I\setminus\bigcup_{(q,q')\in K}\ooint{q,q'}$ mempunyai paling banyak
dua titik (titik-titik ujung $I$, jika ada), sehingga terabaikan. Jika
kita menuliskan $S$ untuk himpunan titik-titik $I$ tempat $f$ tidak
terdiferensialkan, maka dari (a)-(c) kita melihat bahwa
$S\cap\ooint{q,q'}$ terabaikan untuk setiap $(q,q')\in K$, sehingga
$S\cap\bigcup_{(q,q')\in K}\ooint{q,q'}$ terabaikan dan $S$ terabaikan.
     
\medskip
     
{\bf (e)} Jadi kita selesai jika $f$ tak-menurun. Untuk $f$ yang
tak-menaik, terapkan hasil di atas pada $-f$, yang terdiferensialkan
tepat di titik-titik yang sama seperti $f$.
}%end of proof of 222A
     
\leader{222B}{Catatan-catatan}\dvro{ Jika}{ {\bf (a)} Saya mencatat bahwa
dalam argumen di atas saya menggunakan rumus-rumus seperti
$\sum_{J\in\Cal I_0}\mu f^*(J)$. Ini karena teorema Vitali tidak menentukan
apakah keluarga $\Cal I_0$ akan berhingga atau tak berhingga. Jumlah
tersebut harus ditafsirkan menurut cara yang ditetapkan dalam 112Bd di
Jilid 1; secara umum, $\sum_{k\in K}a_k$, dengan $K$ suatu himpunan
sembarang dan setiap $a_k\ge 0$, dimaksudkan sebagai
$\sup_{L\subseteq K\text{ berhingga}}\sum_{k\in L}a_k$, dengan konvensi
bahwa $\sum_{k\in\emptyset}a_k=0$. Sekarang, dalam konteks ini, jika}
$(X,\Sigma,\mu)$ merupakan ruang ukur, $K$ suatu himpunan terhitung,
dan $\family{k}{K}{E_k}$ suatu keluarga dalam $\Sigma$,
     
\Centerline{$\mu(\bigcup_{k\in K}E_k)\le\sum_{k\in K}\mu E_k$,}
     
\noindent dengan kesamaan jika $\family{k}{K}{E_k}$ saling lepas.
\prooflet{\Prf\ Jika $K=\emptyset$, hal ini sepele. Jika tidak, misalkan
$n\mapsto k_n:\Bbb N\to K$ suatu surjeksi, dan tetapkan
     
\Centerline{$K_n=\{k_i:i\le n\}$, \quad$G_n=\bigcup_{i\le n}E_{k_i}
=\bigcup_{k\in K_n}E_k$}
     
\noindent untuk setiap $n\in\Bbb N$. Maka $\sequencen{G_n}$ merupakan
barisan tak-menurun dengan gabungan $E=\bigcup_{k\in K}E_k$, sehingga
     
\Centerline{$\mu E=\lim_{n\to\infty}\mu G_n=\sup_{n\in\Bbb N}\mu G_n$;}
     
\noindent dan
     
\Centerline{$\mu G_n\le\sum_{k\in K_n}\mu E_k\le\sum_{k\in K}\mu E_k$}
     
\noindent untuk setiap $n$, sehingga
$\mu E\le\sum_{k\in K}\mu E_k$. Jika $E_k$ saling lepas, maka $\mu G_n$
tepat sama dengan $\sum_{k\in K_n}\mu E_k$ untuk setiap $n$; tetapi
karena $\sequencen{K_n}$ merupakan barisan tak-menurun dari himpunan-himpunan
dengan gabungan $K$, setiap subhimpunan berhingga dari $K$ termasuk dalam
suatu $K_n$, dan
     
\Centerline{$\sum_{k\in K}\mu E_k=\sup_{n\in\Bbb N}\sum_{k\in K_n}
\mu E_k=\sup_{n\in\Bbb N}\mu G_n=\mu E$,}
     
\noindent sebagaimana diperlukan.\ \Qed}
     
\cmmnt{
\spheader 222Bb Sebagian pembaca akan lebih menyukai pengindeksan ulang
himpunan-himpunan secara teratur, sehingga semua jumlah yang perlu mereka
perhatikan berbentuk $\sum_{i=0}^n$ atau $\sum_{i=0}^{\infty}$. Pada
hakikatnya, itulah yang saya lakukan dalam Jilid 1, dalam bukti
114Da/115Da, ketika menunjukkan bahwa ukuran luar Lebesgue memang
merupakan ukuran luar. Kerugian prosedur ini dalam konteks 222A ialah
bahwa kita harus terus-menerus memeriksa bahwa tidak menjadi soal apakah
pada suatu saat tertentu kita mempunyai jumlah berhingga atau tak
berhingga. Saya percaya bahwa teknik yang diuraikan di sini layak
dipelajari dengan sungguh-sungguh, karena sangat sering kita hendak
meninjau gabungan himpunan-himpunan yang diindeks oleh himpunan selain
$\Bbb N$ dan $\{0,\ldots,n\}$.
     
\spheader 222Bc Tentu saja argumen di atas dapat dipersingkat jika Anda
mengetahui sedikit lebih banyak tentang himpunan terhitung daripada yang
telah saya nyatakan secara eksplisit sejauh ini. Tetapi perhatikan bahwa
nilai yang diberikan kepada $\sum_{k\in K}a_k$ tidak boleh bergantung
pada enumerasi $\sequencen{k_n}$ yang kita pilih.
}%end of comment
     
     
\leader{222C}{Lema} Andaikan bahwa $a\le b$ dalam $\Bbb R$, dan bahwa
$F:[a,b]\to\Bbb R$ merupakan fungsi tak-menurun. Maka $\int_a^bF'$ ada
dan paling besar $F(b)-F(a)$.
     
\medskip
     
\def\tmpins{}
     
\noindent{\bf Catatan}\cmmnt{ Saya membahas integrasi atas subhimpunan
secara panjang lebar dalam \S131 dan \S214. Untuk subhimpunan terukur,
yang sudah cukup bagi keperluan kita dalam bab ini, kita mempunyai
deskripsi sederhana: jika $(X,\Sigma,\mu)$ suatu ruang ukur, $E\in\Sigma$
dan $f$ suatu fungsi bernilai real, maka $\int_Ef=\int\tilde f$ jika
integral yang terakhir ada, dengan
$\dom\tilde f=(E\cap\dom f)\cup(X\setminus E)$ dan $\tilde f(x)=f(x)$
jika $x\in E\cap\dom f$, serta $0$ jika $x\in X\setminus E$ (terapkan
131Fa pada $\tilde f$). Segera diperoleh bahwa jika sekarang $F\in\Sigma$
dan $F\subseteq E$, maka $\int_Ff=\int_Ef\times\chi F$.
     
} Saya menuliskan $\int_a^xf$ untuk mengartikan
$\int_{\coint{a,x}}f$\cmmnt{, yang (karena $\coint{a,x}$ terukur) dapat
ditangani sebagaimana diuraikan di atas}.
\cmmnt{Perhatikan bahwa selama kita menangani ukuran Lebesgue, sehingga
$[a,x]\setminus\ooint{a,x}=\{a,x\}$ terabaikan, kita tidak perlu
membedakan $\int_{[a,x]}$, $\int_{\ooint{a,x}}$,
$\int_{\coint{a,x}}$, dan $\int_{\ocint{a,x}}$; untuk ukuran lain pada
$\Bbb R$ kita mungkin harus lebih berhati-hati. Saya menggunakan interval
setengah terbuka agar jelas bahwa
$\int_a^xf+\int_x^yf=\int_a^yf$ jika $a\le x\le y$, karena
     
\Centerline{$f\times\chi\coint{a,y}
=f\times\chi\coint{a,x}\,+\,f\times\chi\coint{x,y}$.}
}%end of comment
     
\proof{{\bf (a)} Hasil ini sepele jika $a=b$; andaikan bahwa $a<b$.
Menurut 222A, $F'$ terdefinisi hampir di mana-mana dalam $[a,b]$.
     
\medskip
     
{\bf (b)} Untuk setiap $n\in\Bbb N$, definisikan fungsi sederhana
$g_n:\coint{a,b}\to\Bbb R$ sebagai berikut. Untuk $0\le k<2^n$, tetapkan
$a_{nk}=a+2^{-n}k(b-a)$, $b_{nk}=a+2^{-n}(k+1)(b-a)$,
$I_{nk}=\coint{a_{nk},b_{nk}}$. Untuk setiap $x\in\coint{a,b}$, ambil
$k<2^n$ sedemikian sehingga $x\in I_{nk}$, dan tetapkan
     
\Centerline{$g_n(x)=\Bover{2^n}{b-a}(F(b_{nk})-F(a_{nk}))$}
     
\noindent untuk $x\in I_{nk}$, sehingga $g_n$ memberikan kemiringan tali
busur grafik $F$ yang ditentukan oleh titik-titik ujung $I_{nk}$. Maka
     
\Centerline{$\int_a^b g_n=\sum_{k=0}^{2^n-1}F(b_{nk})-F(a_{nk})
=F(b)-F(a)$.}
     
\medskip
     
{\bf (c)} Di sisi lain, jika kita tetapkan
     
\Centerline{$C=\{x:x\in\ooint{a,b},\,F'(x)$ ada$\}$,}
     
\noindent maka $[a,b]\setminus C$ terabaikan, menurut 222A, dan
$F'(x)=\lim_{n\to\infty}g_n(x)$ untuk setiap $x\in C$. \Prf\ Misalkan
$\epsilon>0$. Maka terdapat $\delta>0$ sedemikian sehingga $x+h\in[a,b]$
dan $|F(x+h)-F(x)-hF'(x)|\le\epsilon|h|$ kapan pun
$|h|\le\delta$. Misalkan $n\in\Bbb N$ sedemikian sehingga
$2^{-n}(b-a)\le\delta$. Misalkan $k<2^n$ sedemikian sehingga
$x\in I_{nk}$. Maka
     
\Centerline{$x-\delta\le a_{nk}\le x<b_{nk}\le x+\delta$,
\quad$g_n(x)=\Bover{2^n}{b-a}(F(b_{nk})-F(a_{nk}))$.}
     
\noindent Sekarang kita mempunyai
     
$$\eqalign{|g_n(x)-F'(x)|
&=|\Bover{2^n}{b-a}(F(b_{nk})-F(a_{nk}))-F'(x)|\cr
&=\Bover{2^n}{b-a}|F(b_{nk})-F(a_{nk})-(b_{nk}-a_{nk})F'(x)|\cr
&\le\Bover{2^n}{b-a}\bigl(|F(b_{nk})-F(x)-(b_{nk}-x)F'(x)|\cr
&\qquad\qquad\qquad\qquad   +|F(x)-F(a_{nk})-(x-a_{nk})F'(x)|\bigr)\cr
&\le \Bover{2^n}{b-a}(\epsilon|b_{nk}-x|+\epsilon|x-a_{nk}|)
=\epsilon.\cr}$$
     
\noindent Dan ini benar kapan pun $2^{-n}(b-a)\le\delta$, yakni untuk semua
$n$ yang cukup besar. Karena $\epsilon$ sembarang,
$F'(x)=\lim_{n\to\infty}g_n(x)$.\ \Qed
     
\medskip
     
{\bf (d)} Jadi $g_n\to F'$ hampir di mana-mana dalam $[a,b]$. Menurut
Lema Fatou,
     
\Centerline{$\int_a^bF'
=\int_a^b\liminf_{n\to\infty}g_n
\le\liminf_{n\to\infty}\int_a^bg_n
=\lim_{n\to\infty}\int_a^b g_n
=F(b)-F(a)$,}
     
\noindent sebagaimana diperlukan.
}%end of proof of 222C
     
\cmmnt{
\medskip
     
\noindent{\bf Catatan} Terdapat suatu generalisasi hasil ini dalam 224I.
}
     
\leader{222D}{Lema} Andaikan bahwa $a<b$ dalam $\Bbb R$, dan bahwa $f$,
$g$ merupakan fungsi bernilai real, keduanya terintegralkan pada $[a,b]$,
sedemikian sehingga $\int_a^xf=\int_a^xg$ untuk setiap $x\in[a,b]$.
Maka $f=g$ hampir di mana-mana dalam $[a,b]$.
     
\proof{ Intinya ialah bahwa
     
\Centerline{$\int_Ef=\int_a^bf\times\chi E
=\int_a^bg\times\chi E=\int_Eg$}
     
\noindent untuk setiap himpunan terukur $E\subseteq\coint{a,b}$.
     
\Prf\ {\bf (i)} Jika $E=\coint{c,d}$ dengan $a\le c\le d\le b$, maka
     
\Centerline{$\int_Ef=\int_a^df-\int_a^cf
=\int_a^dg-\int_a^cg=\int_Eg$.}
     
\medskip
     
\quad{\bf (ii)} Jika $E=\coint{a,b}\cap G$ untuk suatu himpunan terbuka
$G\subseteq\Bbb R$, maka untuk setiap $n\in\Bbb N$ tetapkan
     
\Centerline{$K_n
=\{k:k\in\Bbb Z,\,|k|\le 4^n,\,\coint{2^{-n}k,2^{-n}(k+1)}\subseteq
G\}$,}
     
\Centerline{$H_n
=\bigcup_{k\in K_n}\coint{2^{-n}k,2^{-n}(k+1)}\cap\coint{a,b}$;}
     
\noindent maka $\sequencen{H_n}$ merupakan barisan tak-menurun dari
himpunan-himpunan terukur dengan gabungan $E$, sehingga
$f\times\chi E=\lim_{n\to\infty}f\times\chi H_n$, dan (menurut Teorema
Konvergensi Terdominasi Lebesgue, karena $|f\times\chi H_n|\le|f|$
hampir di mana-mana untuk setiap $n$, dan $|f|$ terintegralkan)
     
\Centerline{$\int_Ef=\lim_{n\to\infty}\int_{H_n}f$.}
     
\noindent Pada saat yang sama, setiap $H_n$ merupakan gabungan berhingga
yang saling lepas dari interval-interval setengah terbuka dalam
$\coint{a,b}$, sehingga
     
\Centerline{$\int_{H_n}f
=\sum_{k\in K_n}\int_{\coint{2^{-n}k,2^{-n}(k+1)}\cap\coint{a,b}}f
=\sum_{k\in K_n}\int_{\coint{2^{-n}k,2^{-n}(k+1)}\cap\coint{a,b}}g
=\int_{H_n}g$,}
     
\noindent dan
     
\Centerline{$\int_Eg=\lim_{n\to\infty}\int_{H_n}g
=\lim_{n\to\infty}\int_{H_n}f=\int_Ef$.}
     
\medskip
     
\quad{\bf (iii)} Untuk himpunan terukur umum $E\subseteq\coint{a,b}$,
kita dapat memilih untuk setiap $n\in\Bbb N$ suatu himpunan terbuka
$G_n\supseteq E$ sedemikian sehingga $\mu G_n\le\mu E+2^{-n}$ (134Fa).
Tetapkan $G_n'=\bigcap_{m\le n}G_m$,
$E_n=\coint{a,b}\cap G'_n$ untuk setiap $n$,
     
\Centerline{$F=\coint{a,b}\cap\bigcap_{n\in\Bbb N}G_n
=\bigcap_{n\in\Bbb N}\coint{a,b}\cap G_n'=\bigcap_{n\in\Bbb N}E_n$.}
     
\noindent Maka $E\subseteq F$ dan
     
\Centerline{$\mu F\le\inf_{n\in\Bbb N}\mu G_n=\mu E$,}
     
\noindent sehingga $F\setminus E$ terabaikan dan
$f\times\chi(F\setminus E)$ bernilai nol hampir di mana-mana; akibatnya
$\int_{F\setminus E}f=0$ dan $\int_Ff=\int_Ef$. Di sisi lain,
     
\Centerline{$f\times\chi F=\lim_{n\to\infty}f\times\chi E_n$,}
     
\noindent sehingga, sekali lagi menurut Teorema Konvergensi Terdominasi
Lebesgue,
     
\Centerline{$\int_Ef=\int_Ff=\lim_{n\to\infty}\int_{E_n}f$.}
     
\noindent Demikian pula,
     
\Centerline{$\int_Eg=\lim_{n\to\infty}\int_{E_n}g$.}
     
\noindent Tetapi menurut bagian (ii) kita mempunyai
$\int_{E_n}g=\int_{E_n}f$ untuk setiap $n$, sehingga
$\int_Eg=\int_Ef$, sebagaimana diperlukan.\ \Qed
     
Menurut 131Hb, $f=g$ hampir di mana-mana dalam $\coint{a,b}$, dan
karena itu hampir di mana-mana dalam $[a,b]$.
}%end of proof of 222D
     
\leader{222E}{Teorema} Andaikan bahwa $a\le b$ dalam $\Bbb R$ dan $f$
merupakan fungsi bernilai real yang terintegralkan pada $[a,b]$. Maka
$F(x)=\int_a^xf$ ada dalam $\Bbb R$ untuk setiap $x\in[a,b]$, dan
turunan $F'(x)$ ada serta sama dengan $f(x)$ untuk hampir setiap
$x\in[a,b]$.
     
\proof{{\bf (a)} Dalam sebagian besar bukti ini (sampai akhir (c) di
bawah) saya mengandaikan bahwa $f$ nonnegatif. Dalam hal ini,
     
     
\Centerline{$F(y)=F(x)+\int_x^yf\ge F(x)$}
     
\noindent kapan pun $a\le x\le y\le b$; jadi $F$ tak-menurun dan karena
itu terdiferensialkan hampir di mana-mana dalam $[a,b]$, menurut 222A.
     
Menurut 222C kita juga mengetahui bahwa $\int_a^xF'$ ada dan kurang dari
atau sama dengan $F(x)-F(a)=F(x)$ untuk setiap $x\in[a,b]$.
     
\medskip
     
{\bf (b)} Sekarang andaikan, sebagai tambahan, bahwa $f$ terbatas;
katakanlah $0\le f(t)\le M$ untuk setiap $t\in\dom f$. Maka $M-f$
terintegralkan pada $[a,b]$; misalkan $G$ integral tak tentunya, sehingga
$G(x)=M(x-a)-F(x)$ untuk setiap $x\in[a,b]$. Dengan menerapkan (a) pada
$M-f$ dan $G$, kita memperoleh $\int_a^xG'\le G(x)$ untuk setiap
$x\in[a,b]$; tetapi tentu saja $G'=M-F'$, sehingga
$M(x-a)-\int_a^xF'\le M(x-a)-F(x)$, yakni $\int_a^xF'\ge F(x)$ untuk
setiap $x\in[a,b]$. Jadi $\int_a^xF'=\int_a^xf$ untuk setiap
$x\in[a,b]$. Sekarang 222D memberi tahu kita bahwa $F'=f$ hampir di
mana-mana dalam $[a,b]$.
     
\medskip
     
{\bf (c)} Jadi kita selesai untuk $f$ yang terbatas dan nonnegatif.
Untuk $f$ yang tak terbatas, misalkan $\sequencen{f_n}$ suatu barisan
tak-menurun dari fungsi sederhana nonnegatif yang konvergen ke $f$
hampir di mana-mana dalam $[a,b]$, dan misalkan $\sequencen{F_n}$
integral-integral tak tentu yang bersesuaian. Maka untuk setiap $n$ dan
setiap $x$, $y$ dengan $a\le x\le y\le b$, kita mempunyai
     
\Centerline{$F(y)-F(x)=\int_x^yf\ge\int_x^yf_n=F_n(y)-F_n(x)$,}
     
\noindent sehingga $F'(x)\ge F_n'(x)$ untuk setiap
$x\in\ooint{a,b}$ tempat keduanya terdefinisi, dan
$F'(x)\ge f_n(x)$ untuk hampir setiap $x\in[a,b]$. Ini benar untuk
setiap $n$, sehingga $F'\ge f$ hampir di mana-mana, dan $F'-f\ge 0$
hampir di mana-mana. Di sisi lain, sebagaimana dicatat dalam (a),
     
\Centerline{$\int_a^bF'\le F(b)-F(a)=\int_a^bf$,}
     
\noindent sehingga $\int_a^bF'-f\le 0$. Akibatnya $F'\eae f$ (yakni,
$F'=f$ hampir di mana-mana dalam $[a,b]$)(122Rd).
     
\medskip
     
{\bf (d)} Ini menyelesaikan bukti untuk $f$ nonnegatif. Untuk $f$
umum, kita dapat menyatakan $f$ sebagai $f_1-f_2$ dengan $f_1$, $f_2$
fungsi terintegralkan nonnegatif; sekarang $F=F_1-F_2$ dengan $F_1$,
$F_2$ integral-integral tak tentu yang bersesuaian, sehingga
$F'\eae F_1'-F_2'\eae f_1-f_2$, dan $F'\eae f$.
}%end of proof of 222E
     
\leader{222F}{Korolari} Andaikan bahwa $f$ sembarang fungsi bernilai real
yang terintegralkan pada $\Bbb R$, dan tetapkan
$F(x)=\int_{-\infty}^xf$ untuk setiap $x\in\Bbb R$. Maka $F'(x)$ ada
dan sama dengan $f(x)$ untuk hampir setiap $x\in\Bbb R$.
     
\proof{ Untuk setiap $n\in\Bbb N$, tetapkan
     
\Centerline{$F_n(x)=\int_{-n}^xf$}
     
\noindent untuk $x\in[-n,n]$. Maka $F_n'(x)=f(x)$ untuk hampir setiap
$x\in[-n,n]$. Tetapi $F(x)=F(-n)+F_n(x)$ untuk setiap $x\in[-n,n]$,
sehingga $F'(x)$ ada dan sama dengan $F'_n(x)$ untuk setiap
$x\in\ooint{-n,n}$ tempat $F'_n(x)$ terdefinisi; dan $F'(x)=f(x)$
untuk hampir setiap $x\in[-n,n]$. Karena $n$ sembarang, $F'\eae f$.
}%end of proof of 222F
     
\leader{222G}{Korolari} Andaikan bahwa $E\subseteq\Bbb R$ suatu himpunan
terukur dan $f$ suatu fungsi bernilai real yang terintegralkan pada $E$.
Tetapkan $F(x)=\int_{E\cap\ooint{-\infty,x}}f$ untuk $x\in\Bbb R$.
Maka $F'(x)=f(x)$ untuk hampir setiap $x\in E$, dan $F'(x)=0$ untuk
hampir setiap $x\in\Bbb R\setminus E$.
     
\proof{ Terapkan 222F pada $\tilde f$, dengan $\tilde f(x)=f(x)$ untuk
$x\in E\cap\dom f$ dan $\tilde f(x)=0$ untuk
$x\in\Bbb R\setminus E$, sehingga $F(x)=\int_{-\infty}^x\tilde f$
untuk setiap $x\in\Bbb R$.
}%end of proof of 222G
     
\leader{222H}{}\cmmnt{ Hasil bahwa $\bover{d}{dx}\int_a^xf=f(x)$ untuk
hampir setiap $x$ memang memuaskan, tetapi tidak menggantikan hasil yang
lebih elementer bahwa kesamaan ini berlaku di setiap titik tempat $f$
kontinu.
     
\medskip
     
\noindent}{\bf Proposisi} Andaikan bahwa $a\le b$ dalam $\Bbb R$ dan
$f$ merupakan fungsi bernilai real yang terintegralkan pada $[a,b]$.
Tetapkan $F(x)=\int_a^xf$ untuk $x\in[a,b]$. Maka $F'(x)$ ada dan sama
dengan $f(x)$ di setiap titik $x\in\dom(f)\cap\ooint{a,b}$ tempat $f$
kontinu.
     
\proof{ Tetapkan $c=f(x)$. Misalkan $\epsilon>0$. Misalkan $\delta>0$
sedemikian sehingga $\delta\le\min(b-x,x-a)$ dan
$|f(t)-c|\le\epsilon$ kapan pun $t\in\dom f$ dan $|t-x|\le\delta$.
Jika $x<y\le x+\delta$, maka
     
\Centerline{$|\Bover{F(y)-F(x)}{y-x}-c|
=\Bover1{y-x}|\int_x^yf-c|
\le\Bover1{y-x}\int_x^y|f-c|
\le\epsilon$.}
     
\noindent Demikian pula, jika $x-\delta\le y<x$,
     
\Centerline{$|\Bover{F(y)-F(x)}{y-x}-f(x)|
=\Bover1{x-y}|\int_y^xf-c|
\le\Bover1{x-y}\int_y^x|f-c|
\le\epsilon$.}
     
\noindent Karena $\epsilon$ sembarang,
     
\Centerline{$F'(x)=\lim_{y\to x}\Bover{F(y)-F(x)}{y-x}=c$,}
     
\noindent sebagaimana diperlukan.
}%end of proof of 222H
     
     
\leader{222I}{Fungsi bernilai kompleks} \cmmnt{Sejauh ini dalam bagian
ini, saya mengambil setiap $f$ bernilai real. Perluasannya ke $f$
bernilai kompleks cukup dilakukan dengan menerapkan hasil-hasil di atas
pada bagian real dan bagian imajiner $f$. Secara khusus, kita memperoleh
yang berikut.

\medskip

}{\bf (a)} Jika $a\le b$ dalam $\Bbb R$ dan $f$ merupakan fungsi
bernilai kompleks yang terintegralkan pada $[a,b]$, maka
$F(x)=\int_a^xf$ terdefinisi dalam $\Bbb C$ untuk setiap $x\in[a,b]$,
dan turunannya $F'(x)$ ada serta sama dengan $f(x)$ untuk hampir setiap
$x\in[a,b]$; selain itu, $F'(x)=f(x)$ kapan pun
$x\in\dom(f)\cap\ooint{a,b}$ dan $f$ kontinu di $x$.
     
\spheader 222Ib Jika $f$ merupakan fungsi bernilai kompleks yang
terintegralkan pada $\Bbb R$, dan $F(x)=\int_{-\infty}^xf$ untuk setiap
$x\in\Bbb R$, maka $F'$ ada dan sama dengan $f$ hampir di mana-mana
dalam $\Bbb R$.
     
\spheader 222Ic Jika $E\subseteq\Bbb R$ suatu himpunan terukur dan $f$
merupakan fungsi bernilai kompleks yang terintegralkan pada $E$, dan
$F(x)=\int_{E\cap\ooint{-\infty,x}}f$ untuk setiap $x\in\Bbb R$, maka
$F'(x)=f(x)$ untuk hampir setiap $x\in E$ dan $F'(x)=0$ untuk hampir
setiap $x\in\Bbb R\setminus E$.
     
\leader{*222J}{Teorema \dvrocolon{Denjoy-Young-Saks}}\cmmnt{ Hasil
berikutnya, menurut rencana saat ini, tidak akan digunakan di bagian mana pun
dalam risalah ini. Namun hasil ini bersifat sentral bagi bagian-bagian
analisis real yang hendak diberi fondasi oleh jilid ini, dan meskipun
argumennya memerlukan kecanggihan tertentu, sebenarnya argumen itu
tidak jauh melampaui teorema Lebesgue 222A. Saya harus memulai dengan
beberapa notasi.
     
\medskip

\noindent}{\bf Definisi} Misalkan $f$ sembarang fungsi bernilai real, dan
$A\subseteq\Bbb R$ domainnya. Tuliskan

\Centerline{$\tilde A^+
=\{x:x\in A$, $\ocint{x,x+\delta}\cap A\ne\emptyset$
untuk setiap $\delta>0\}$,}

\Centerline{$\tilde A^-
=\{x:x\in A$, $\coint{x-\delta,x}\cap A\ne\emptyset$
untuk setiap $\delta>0\}$.}
     
\noindent Tetapkan

\Centerline{$\DiniD^+(x)
=\limsup_{y\in A,y\downarrow x}\Bover{f(y)-f(x)}{y-x}
=\inf_{\delta>0}\sup_{y\in A,x<y\le x+\delta}\Bover{f(y)-f(x)}{y-x}$,}

\Centerline{$\Dinid^+(x)
=\liminf_{y\in A,y\downarrow x}\Bover{f(y)-f(x)}{y-x}
=\sup_{\delta>0}\inf_{y\in A,x<y\le x+\delta}\Bover{f(y)-f(x)}{y-x}$}

\noindent untuk $x\in\tilde A^+$, dan

\Centerline{$\DiniD^-(x)
=\limsup_{y\in A,y\uparrow x}\Bover{f(y)-f(x)}{y-x}
=\inf_{\delta>0}\sup_{y\in A,x-\delta\le y<x}\Bover{f(y)-f(x)}{y-x}$,}

\Centerline{$\Dinid^-(x)
=\liminf_{y\in A,y\uparrow x}\Bover{f(y)-f(x)}{y-x}
=\sup_{\delta>0}\inf_{y\in A,x-\delta\le y<x}\Bover{f(y)-f(x)}{y-x}$}

\noindent untuk $x\in\tilde A^-$, semuanya terdefinisi dalam
$[-\infty,\infty]$. (Inilah keempat {\bf turunan Dini} dari
$f$.\cmmnt{ Anda juga akan menjumpai $D^+$, $d^+$, $D^-$, $d^-$ yang
digunakan sebagai pengganti $\DiniD^+$, $\Dinid^+$, $\DiniD^-$, dan
$\Dinid^-$.})

Perhatikan bahwa tentu saja kita mempunyai
$(\Dinid^+f)(x)\le(\DiniD^+f)(x)$ untuk setiap $x\in\tilde A^+$,
sedangkan $(\Dinid^-f)(x)\le(\DiniD^-f)(x)$ untuk setiap
$x\in\tilde A^-$. Turunan biasa $f'(x)$ terdefinisi dan sama dengan
$c\in\Bbb R$ jika dan hanya jika ($\alpha$) $x$ termasuk dalam suatu
interval terbuka yang termuat dalam $A$ ($\beta$)
$(\DiniD^+f)(x)=(\Dinid^+f)(x)=(\DiniD^-f)(x)
=(\Dinid^-f)(x)=c$.

\leader{*222K}{Lema} Misalkan $A$ sembarang subhimpunan $\Bbb R$, dan
definisikan $\tilde A^+$ dan $\tilde A^-$ seperti dalam 222J. Maka
$A\setminus\tilde A^+$ dan $A\setminus\tilde A^-$ terhitung, dan karena
itu terabaikan.

\proof{ Kita mempunyai

\Centerline{$A\setminus\tilde A^+
=\bigcup_{q\in\Bbb Q}\{x:x\in A$, $x<q$, $A\cap\ocint{x,q}=\emptyset\}$.}

\noindent Tetapi untuk setiap $q\in\Bbb Q$ dapat ada paling banyak satu
$x\in A$ sedemikian sehingga $x<q$ dan $\ocint{x,q}$ tidak bertemu $A$,
sehingga $A\setminus\tilde A^+$ merupakan gabungan terhitung dari
himpunan-himpunan berhingga dan bersifat terhitung. Demikian pula,

\Centerline{$A\setminus\tilde A^-
=\bigcup_{q\in\Bbb Q}\{x:x\in A$, $q<x$, $A\cap\coint{q,x}=\emptyset\}$}

\noindent terhitung.
}%end of proof of 222K

\leader{*222L}{Teorema} Misalkan $f$ sembarang fungsi bernilai real, dan $A$
domainnya. Maka untuk hampir setiap $x\in A$

\qquad{\it entah} keempat turunan Dini dari $f$ di $x$ terdefinisi,
berhingga, dan sama

\qquad{\it atau} $(\DiniD^+f)(x)=(\Dinid^-f)(x)$ berhingga,
$(\Dinid^+f)(x)=-\infty$ dan $(\DiniD^-f)(x)=\infty$

\qquad{\it atau} $(\Dinid^+f)(x)=(\DiniD^-f)(x)$ berhingga,
$(\DiniD^+f)(x)=\infty$ dan $(\Dinid^-f)(x)=-\infty$

\qquad{\it atau} $(\DiniD^+f)(x)=(\DiniD^-f)(x)=\infty$ dan
$(\Dinid^+f)(x)=(\Dinid^-f)(x)=-\infty$.

\proof{{\bf (a)} Tetapkan $\tilde A=\tilde A^+\cap\tilde A^-$, dengan
mendefinisikan $\tilde A^+$ dan $\tilde A^-$ seperti dalam 222J,
sehingga komplemen $\tilde A$ dalam $A$ terhitung dan keempat turunan Dini
terdefinisi pada $\tilde A$. Untuk $n\in\Bbb N$, $q\in\Bbb Q$
tetapkan

\Centerline{$E_{qn}=\{x:x\in\tilde A$, $x<q$, $f(y)\ge f(x)-n(y-x)\}$
untuk setiap $y\in A\cap[x,q]\}$.}

\noindent Perhatikan bahwa

\Centerline{$\bigcup_{n\in\Bbb N,q\in\Bbb Q}E_{qn}=\{x:x\in\tilde A$,
$(\Dinid^+f)(x)>-\infty\}$.}

\noindent Untuk $q\in\Bbb Q$, $n\in\Bbb N$ sedemikian sehingga
$E_{qn}$ tidak kosong, tetapkan
$\beta_{qn}=\sup E_{qn}\in\ocint{-\infty,q}$,
$\alpha_{qn}=\inf E_{qn}\in[-\infty,\beta_{qn}]$, dan untuk
$x\in\ooint{\alpha_{qn},\beta_{qn}}$ tetapkan
$g_{qn}(x)=\inf\{f(y)+ny:y\in A\cap[x,q]\}$. Perhatikan bahwa jika
$x\in E_{qn}\setminus\{\alpha_{qn},\beta_{qn}\}$, maka
$g_{qn}(x)=f(x)+nx$ berhingga; juga $g_{qn}$ monoton, sehingga berhingga di
setiap titik dalam $\ooint{\alpha_{qn},\beta_{qn}}$, dan tentu saja
$g_{qn}(x)\le f(x)+nx$ untuk setiap
$x\in A\cap\ooint{\alpha_{qn},\beta_{qn}}$.

Menurut 222A, hampir setiap titik $\ooint{\alpha_{qn},\beta_{qn}}$
termasuk dalam $F_{qn}=\dom g'_{qn}$; khususnya,
$E_{qn}\setminus F_{qn}$ terabaikan. Tetapkan
$h_{qn}(x)=g_{qn}(x)-nx$ untuk
$x\in\ooint{\alpha_{qn},\beta_{qn}}$; maka $h_{qn}$ terdiferensialkan di
setiap titik $F_{qn}$. Sekarang jika $x\in E_{qn}\cap F_{qn}$, kita
mempunyai $h_{qn}(x)=f(x)$, sedangkan $h_{qn}(x)\le f(x)$ untuk
$x\in A\cap\ooint{\alpha_{qn},\beta_{qn}}$; akibatnya

$$\eqalign{(\Dinid^+f)(x)
&=\sup_{\delta>0}\inf_{y\in A\cap\ocint{x,x+\delta}}
  \Bover{f(y)-f(x)}{y-x}\cr
&\ge\sup_{\delta>0}\inf_{y\in A\cap\ocint{x,x+\delta}}
  \Bover{h_{qn}(y)-h_{qn}(x)}{y-x}\cr
&\ge\sup_{0<\delta<\beta_{qn}-x}\inf_{y\in\ocint{x,x+\delta}}
  \Bover{h_{qn}(y)-h_{qn}(x)}{y-x}
=h_{qn}'(x),\cr}$$
  
\noindent dan demikian pula
$(\DiniD^-f)(x)\le h'_{qn}(x)$. Di sisi lain, jika
$x\in\tilde E_{qn}^+$, maka

$$\eqalign{(\Dinid^+f)(x)
&=\sup_{\delta>0}\inf_{y\in A\cap\ocint{x,x+\delta}}
    \Bover{f(y)-f(x)}{y-x}\cr
&\le\sup_{\delta>0}\inf_{y\in E_{qn}\cap\ocint{x,x+\delta}}
    \Bover{f(y)-f(x)}{y-x}\cr
&=\sup_{\delta>0}\inf_{y\in E_{qn}\cap\ocint{x,x+\delta}}
    \Bover{h_{qn}(y)-h_{qn}(x)}{y-x}\cr
&\le\inf_{\delta>0}\sup_{y\in E_{qn}\cap\ocint{x,x+\delta}}
    \Bover{h_{qn}(y)-h_{qn}(x)}{y-x}\cr
&\le\inf_{0<\delta<\beta_{qn}-x}\sup_{y\in\ocint{x,x+\delta}}
    \Bover{h_{qn}(y)-h_{qn}(x)}{y-x}
=h'_{qn}(x).\cr}$$

\noindent Dengan menggabungkan semua ini, kita melihat bahwa jika
$x\in F_{qn}\cap\tilde E_{qn}^+$, maka
$(\Dinid^+f)(x)=h'_{qn}(x)\ge(\DiniD^-f)(x)$.

Dengan konvensi $F_{qn}=\emptyset$ jika $E_{qn}$
kosong, kalimat terakhir benar untuk semua $q\in\Bbb Q$ dan
$n\in\Bbb N$, sedangkan $A\setminus\tilde A$ dan
$\bigcup_{q\in\Bbb Q,n\in\Bbb N}E_{qn}\setminus(F_{qn}\cap\tilde E_{qn}^+)$
terabaikan, dan $(\Dinid^+f)(x)=-\infty$ untuk setiap
$x\in\tilde A\setminus\bigcup_{q\in\Bbb Q,n\in\Bbb N}E_{qn}$. Jadi kita
melihat bahwa, untuk hampir setiap $x\in A$, entah
$(\Dinid^+f)(x)=-\infty$ atau
$\infty>(\Dinid^+f)(x)\ge(\DiniD^-f)(x)$.

\medskip

{\bf (b)} Dengan mencerminkan argumen di atas dari kiri ke kanan atau
dari atas ke bawah, kita melihat bahwa, untuk hampir setiap $x\in A$,

\Centerline{entah $(\Dinid^-f)(x)=-\infty$ atau
$\infty>(\Dinid^-f)(x)\ge(\DiniD^+f)(x)$,}

\Centerline{entah $(\DiniD^+f)(x)=\infty$ atau
$-\infty<(\DiniD^+f)(x)\le(\Dinid^-f)(x)$,}

\Centerline{entah $(\DiniD^-f)(x)=\infty$ atau
$-\infty<(\DiniD^-f)(x)\le(\Dinid^+f)(x)$,}

\noindent dan juga

\Centerline{$(\Dinid^+f)(x)\le(\DiniD^+f)(x)$,
\quad$(\Dinid^-f)(x)\le(\DiniD^-f)(x)$.}

\noindent Oleh karena itu, untuk $x$ seperti itu,

\Centerline{$(\Dinid^+f)(x)>-\infty
\Longrightarrow (\DiniD^-f)(x)\le(\Dinid^+f)(x)<\infty
\Longrightarrow (\Dinid^+f)(x)=(\DiniD^-f)(x)\in\Bbb R$,}

\noindent dan demikian pula

\Centerline{$(\Dinid^-f)(x)>-\infty
\Longrightarrow(\Dinid^-f)(x)=(\DiniD^+f)(x)\in\Bbb R$,}

\Centerline{$(\DiniD^+f)(x)<\infty
\Longrightarrow(\DiniD^+f)(x)=(\Dinid^-f)(x)\in\Bbb R$,}

\Centerline{$(\DiniD^-f)(x)<\infty
\Longrightarrow(\DiniD^-f)(x)=(\Dinid^+f)(x)\in\Bbb R$.}

\noindent Jadi kita mempunyai

\Centerline{entah $(\Dinid^+f)(x)=(\DiniD^-f)(x)$ berhingga
atau $(\Dinid^+f)(x)=-\infty$ dan
$(\DiniD^-f)(x)=\infty$,}

\Centerline{entah $(\Dinid^-f)(x)=(\DiniD^+f)(x)$ berhingga
atau $(\Dinid^-f)(x)=-\infty$ dan
$(\DiniD^+f)(x)=\infty$.}

\noindent Kedua dikotomi ini menghasilkan empat kemungkinan; dan karena

\Centerline{$(\Dinid^+f)(x)=(\DiniD^-f)(x)$ berhingga,
\quad$(\Dinid^-f)(x)=(\DiniD^+f)(x)$ berhingga}

\noindent dapat benar bersamaan hanya ketika keempat turunan Dini sama dan
berhingga, kita memperoleh empat kasus yang dicantumkan dalam pernyataan
teorema.
}%end of proof of 222L


\exercises{
\leader{222X}{Latihan dasar $\pmb{>}$(a)}
%\sqheader 222Xa
Misalkan $F:[0,1]\to[0,1]$ fungsi Cantor (134H).
Tunjukkan bahwa $\int_0^1F'=0<F(1)-F(0)$.
%222C
     
\sqheader 222Xb Andaikan bahwa $a<b$ dalam $\Bbb R$ dan $h$ suatu
fungsi bernilai real sedemikian sehingga $\int_x^yh$ ada dan nonnegatif
kapan pun $a\le x\le y\le b$. Tunjukkan bahwa $h\ge 0$ hampir di
mana-mana dalam $[a,b]$.
%222D
     
\sqheader 222Xc Andaikan bahwa $a<b$ dalam $\Bbb R$ dan $f$, $g$
merupakan fungsi bernilai kompleks yang terintegralkan pada $[a,b]$
sedemikian sehingga $\int_a^xf=\int_a^xg$ untuk setiap $x\in[a,b]$.
Tunjukkan bahwa $f=g$ hampir di mana-mana dalam $[a,b]$.
%222I
     
\sqheader 222Xd Andaikan bahwa $a<b$ dalam $\Bbb R$ dan $f$ suatu
fungsi bernilai real yang terintegralkan pada $[a,b]$. Tunjukkan bahwa
integral tak tentu $x\mapsto\int_a^xf$ kontinu.
%222+
     
\leader{222Y}{Latihan lanjutan (a)}
%\spheader 222Ya
Misalkan $\sequencen{F_n}$ suatu barisan fungsi nonnegatif dan tak-menurun
pada $[0,1]$ sedemikian sehingga
$F(x)=\sum_{n=0}^{\infty}F_n(x)$ berhingga untuk setiap $x\in[0,1]$.
Tunjukkan bahwa $\sum_{n=0}^{\infty}F'_n(x)=F'(x)$ untuk hampir setiap
$x\in[0,1]$. ({\it Petunjuk\/}: ambil $\sequence{k}{n_k}$ sedemikian
sehingga $\sum_{k=0}^{\infty}F(1)-G_k(1)<\infty$, dengan
$G_k=\sum_{j=0}^{n_k}F_j$, dan tetapkan
$H(x)=\sum_{k=0}^{\infty}F(x)-G_k(x)$. Perhatikan bahwa
$\sum_{k=0}^{\infty}F'(x)-G'_k(x)\le H'(x)$ kapan pun semua turunannya
terdefinisi, sehingga $F'=\lim_{k\to\infty}G'_k$ hampir di mana-mana.)
%222A
     
\spheader 222Yb Misalkan $F:[0,1]\to\Bbb R$ suatu fungsi kontinu
tak-menurun. (i) Tunjukkan bahwa jika $c\in\Bbb R$, maka
$C=\{(x,y):x,\,y\in[0,1],\,F(y)-F(x)=c\}$ terhubung.
({\it Petunjuk\/}: Suatu himpunan $A\subseteq\BbbR^r$ disebut
{\bf terhubung} jika tidak ada surjeksi kontinu $h:A\to\{0,1\}$.
Tunjukkan bahwa jika $h:C\to\{0,1\}$ kontinu, maka fungsi itu berbentuk
$(x,y)\mapsto h_1(x)$ untuk suatu fungsi kontinu $h_1$.) (ii) Sekarang
andaikan bahwa $F(0)=0$, $F(1)=1$ dan bahwa $G:[0,1]\to[0,1]$ merupakan
fungsi kontinu tak-menurun kedua dengan $G(0)=0$, $G(1)=1$. Tunjukkan
bahwa untuk setiap $n\ge 1$ terdapat $x$, $y\in[0,1]$ sedemikian sehingga
$F(y)-F(x)=G(y)-G(x)=\bover1n$.
%222+
     
\spheader 222Yc Misalkan $f$, $g$ fungsi terintegralkan nonnegatif pada
$\Bbb R$, dan $n\ge 1$. Tunjukkan bahwa terdapat $u<v$ dalam
$[-\infty,\infty]$ sedemikian sehingga
$\int_u^vf=\bover1n\int f$ dan $\int_u^vg=\bover1n\int g$.
%222Yb 222+
     
\spheader 222Yd Misalkan $f:\Bbb R\to\Bbb R$ terukur. Tunjukkan bahwa
$H=\dom f'$ merupakan himpunan terukur dan bahwa $f'$ merupakan fungsi
terukur.
%222+

\spheader 222Ye Bangun suatu fungsi terukur Borel
$f:[0,1]\to\{-1,0,1\}$ sedemikian sehingga masing-masing dari empat
kemungkinan yang diuraikan dalam Teorema 222L terjadi pada suatu himpunan
berukuran $\bover14$.
%222L
}%end of exercises
     
\endnotes{
\Notesheader{222} Saya menyerahkan fakta mendasar bahwa suatu integral
tak tentu $x\mapsto\int_a^xf$ selalu kontinu sebagai latihan (222Xd);
sebenarnya fakta ini tidak diperlukan dalam bagian ini, dan suatu hasil yang
jauh lebih kuat diberikan dalam 225E. Masih banyak pula yang harus
dikatakan tentang fungsi monoton, yang akan saya bahas kembali dalam
\S224. Yang kita perlukan di sini ialah fakta bahwa fungsi-fungsi itu
terdiferensialkan hampir di mana-mana (222A), yang saya buktikan dengan
menerapkan teorema Vitali tiga kali, satu kali dalam bagian (b) dari
bukti dan dua kali dalam bagian (c). Setelah ini, argumen-argumen
222C-222E membentuk suatu rangkaian latihan yang bagus tentang gagasan
sentral Jilid 1, dengan menggunakan konsep integrasi atas suatu
subhimpunan (terukur), Lema Fatou (bagian (d) dari bukti 222C), Teorema
Konvergensi Terdominasi Lebesgue (bagian (ii) dan (iii) dari bukti 222D),
dan hampiran himpunan terukur Lebesgue oleh himpunan terbuka (bagian
(iii) dari bukti 222D). Tentu saja, mengetahui bahwa
$\bover{d}{dx}\int_a^xf=f(x)$ hampir di mana-mana sama sekali tidak sama
dengan mengetahui bahwa hal ini berlaku untuk suatu $x$ tertentu, dan
ketika hendak mendiferensialkan suatu integral tak tentu tertentu kita
biasanya terlebih dahulu menggunakan 222H; inti 222E ialah bahwa hasil
itu berlaku pada fungsi yang sangat tidak kontinu, yang sama sekali
tidak dapat ditangani oleh metode yang lebih elementer.

Teorema Denjoy-Young-Saks (222L) merupakan salah satu titik awal suatu
teori subur tentang fenomena `tipikal' dalam analisis real. Mudah untuk
membangun suatu fungsi $f$ dengan sembarang himpunan nilai yang ditentukan
sebelumnya untuk $(\DiniD^+f)(0)$, $(\Dinid^+f)(0)$,
$(\DiniD^-f)(0)$, dan $(\Dinid^-f)(0)$ (tentu saja dengan tunduk pada
syarat $(\Dinid^+f)(0)\le(\DiniD^+f)(0)$ dan
$(\Dinid^-f)(0)\le(\DiniD^-f)(0)$). Tetapi 222L memberi tahu kita bahwa
kombinasi seperti $(\Dinid^-f)(x)=(\Dinid^+f)(x)=\infty$ (yang dapat
kita sebut `$f'(x)=\infty$') hanya dapat terjadi pada himpunan
terabaikan. Empat kemungkinan dalam 222L yang mudah diwujudkan (lihat
222Ye) merupakan satu-satunya yang dapat muncul di titik-titik yang
`tipikal' bagi fungsi yang diberikan, dari sudut pandang ukuran Lebesgue.
Untuk fungsi monoton, 222A memberi tahu kita lebih banyak: pada
titik-titik yang `tipikal' bagi suatu fungsi monoton, fungsi tersebut
benar-benar terdiferensialkan. Dalam bagian berikutnya kita akan melihat
beberapa cara lain untuk menghasilkan himpunan terabaikan dan
koterabaikan dari suatu himpunan atau fungsi yang diberikan, yang
membawa pada penyempurnaan lebih lanjut atas gagasan ini.
}%end of comment
     
\discrpage
     
